% =====================================================================
\part{O argumento}
% =====================================================================

\chapter{A tese em três tamanhos}

Toda banca pede a mesma coisa três vezes, em três tamanhos. Na abertura, quer
uma frase. No meio da apresentação, quer o parágrafo. Na arguição, quer o
raciocínio inteiro. Tenha os três prontos e distintos, e não a mesma coisa
falada mais devagar.

\section{Uma frase}

\begin{quote}\itshape\color{cortinta}
Goiás não abriu uma fronteira nova entre 1985 e 2024: reorganizou a que já
tinha, empurrando lavoura, pastagem e rebanho cerca de setenta quilômetros ao
norte, sem que uma metade do estado tenha empurrado a outra.
\end{quote}

Repare no que a frase faz. Ela afirma um fato robusto (o deslocamento, com
magnitude), nega uma explicação sedutora (o empurrão) e não promete a
explicação alternativa como estabelecida. Essa é exatamente a distribuição de
peso que a evidência sustenta, e é por isso que ela é segura de dizer.

\section{Um parágrafo}

\begin{quote}\color{cortinta}
Em quarenta anos Goiás perdeu 5,8 milhões de hectares de vegetação natural,
passando de 51,9\% a 34,9\% do território, enquanto a agricultura se
multiplicou por 4,8 e a pastagem terminou praticamente onde começou em área,
depois de subir até um pico em 2003 e recuar. O saldo esconde o movimento: os
centros de massa da pastagem, do rebanho e da lavoura deslocaram-se de 65 a 78
quilômetros ao norte, e o fizeram por caminhos opostos. No Sul, a lavoura tomou
o lugar do pasto já aberto; no Norte, o pasto avançou sobre o Cerrado. A
hipótese natural, a de que o Sul tenha empurrado o Norte, foi submetida a
trinta e seis recortes e não apresentou a assinatura que exigiria. O que fica
no lugar dela é uma resposta coordenada a um estímulo comum, distribuída ao
longo do eixo Sul--Norte, como evidência corroborante e não estabelecida. E a
desaceleração recente do Sul não vem de mercado fraco, porque a demanda estava
no pico, mas é compatível com o esgotamento da terra ainda exposta: à medida
que uma unidade se esgota, caem juntos o estoque que lhe resta e a taxa com que
ele é convertido.
\end{quote}

\section{O raciocínio inteiro}

O trabalho é uma cadeia de quatro perguntas em que cada resposta obrigou a
seguinte. Sabê-la nesta ordem é o que permite responder a qualquer pergunta
sem perder o fio.

\vspace{6pt}
\textbf{Primeira: o padrão existe?} A medida é o centro de massa de cada uso da
terra, acompanhado ano a ano. Tudo marchou ao norte, e a marcha sobrevive a
cinco verificações independentes: refeita pixel a pixel sem malha nenhuma,
aberta por formação vegetal, com barra de erro por reamostragem, com
reamostragem de blocos espaciais que não trata as vizinhas como independentes,
e contra dois controles negativos que não acompanham (a área urbana anda para o
\emph{sul}; o leite sobe menos da metade do que sobe o rebanho de corte).
\onde{§4.2}

\textbf{Segunda: que conversão a moveu?} A marcha é o saldo de duas conversões
que não estão no mesmo lugar. No Sul, pastagem vira lavoura; no Norte,
vegetação natural vira pastagem. E a tentativa de explicar \emph{quando} um
pasto é convertido falhou: a região explica 0,5\% da variação da idade, o
sistema produtivo não sustenta e o crédito só aponta na direção contrária à
esperada. A decisão acontece abaixo da escala que dados municipais alcançam.
\onde{§4.3}

\textbf{Terceira: uma região empurrou a outra?} Duas assinaturas foram
procuradas, precedência no tempo e direção no espaço. Nenhuma apareceu. No
lugar do empurrão, o desenho encontrou o maior coeficiente da frente inteira,
que é a \emph{substituição local}: dentro do mesmo município, quando a lavoura
cresce, o pasto recua. Testou-se então a alternativa, um mesmo choque
macroeconômico incidindo com força desigual ao longo do eixo Sul--Norte. O
sinal tem a direção prevista e passa em toda a bateria de especificidade, mas
não cruza o corte de significância sob a inferência apropriada, e a
característica do eixo que faz a diferença não foi isolada. \onde{§4.4}

\textbf{Quarta: por que o Sul freou?} Não foi falta de comprador, pois câmbio,
preço da soja e área de soja plantada estavam todos em alta no período. O Sul
continuou querendo terra --- a conversão de pasto em lavoura acelerou 51\% lá
--- e parou de tirá-la do Cerrado. A evidência do teto de oferta se monta por
eliminação (não é demanda, não é proteção integral) e por uma regularidade
medida: a taxa de conversão cai à medida que a unidade se esgota, resultado que
sobrevive a dezesseis tratamentos diferentes da mesma variável. \onde{§4.5}

\begin{nabanca}[A transição entre as frentes, dita em voz alta]
``A primeira frente me deu um padrão robusto e nenhuma explicação. A segunda me
deu duas lógicas simultâneas e me tirou a explicação regional. A terceira testou
a hipótese que eu mais queria que fosse verdadeira, e ela não passou. A quarta
explica por que a fronteira mudou de endereço em vez de parar.''
\end{nabanca}

\chapter{Por que os nulos são o miolo, e não o problema}

Vários dos resultados centrais deste trabalho são resultados nulos. Isso pode
ser lido de duas maneiras pela banca, e a diferença entre elas depende
inteiramente de como você as apresenta.

\textbf{A leitura ruim} é ``ele não achou nada''. \textbf{A leitura correta} é
que o trabalho testou formalmente a explicação mais atraente disponível, ela
não passou, e a alternativa que ficou no lugar recomenda uma política oposta.
Um nulo com poder declarado e mecanismo alternativo medido é um resultado; um
nulo sem isso é um silêncio.

Por isso, três regras governam toda menção a um nulo neste trabalho.

\vspace{4pt}
\textbf{Regra 1 --- ``não encontrei'' nunca vira ``não há''.} A afirmação
sustentável é que a assinatura testada não aparece nos dados examinados, com o
poder disponível. Ela não é uma refutação do fenômeno, que exigiria desenho e
escala distintos.

\textbf{Regra 2 --- todo nulo declara o seu poder.} O teste temporal do empurrão
descarta um empurrão forte com poder aproximado de 93\%, e não descarta um
moderado, cujo poder é de aproximadamente 48\%. Dizer isso antes de a banca
perguntar é o que separa um nulo defendido de um nulo confessado.

\textbf{Regra 3 --- um nulo mal enunciado afirma mais do que a evidência
sustenta.} ``O deslocamento indireto de uso da terra foi refutado'' é uma frase
proibida neste trabalho, e o Capítulo~\ref{cap:frases} traz a lista completa
delas com as substitutas.

\begin{ancora}[O que os nulos deste trabalho descartam e o que não descartam]
\textbf{Descartado com margem confortável:} um empurrão local visível na série e
um empurrão temporal forte. \\
\textbf{Fora do alcance do teste:} deslocamento de longo alcance que pule os
vizinhos, deslocamento de defasagem longa, efeito temporal moderado que o poder
não alcança, e qualquer canal que atravesse a divisa do estado.
\end{ancora}

\chapter{O mapa do documento}

Saber onde cada coisa está no seu próprio texto é metade da tranquilidade
durante a arguição. O documento tem 115 páginas assim distribuídas.

\vspace{6pt}
\begin{longtable}{@{}p{1.9cm}p{3.2cm}p{9.6cm}@{}}
\toprule
\textbf{Páginas} & \textbf{Parte} & \textbf{O que a banca procura ali} \\
\midrule
\endfirsthead
\toprule
\textbf{Páginas} & \textbf{Parte} & \textbf{O que a banca procura ali} \\
\midrule
\endhead
12--13 & Memorial & Sua trajetória e a declaração sobre o uso de IA \\
14--16 & Cap. 1 --- Introdução & A pergunta, os quatro objetivos específicos, a justificativa \\
17--24 & Cap. 2 --- Referencial & Os arcabouços mobilizados e a lista do que ainda falta ler (§2.10) \\
25--49 & Cap. 3 --- Metodologia & Fontes, AMC, painel, periodização, instrumental, as sete decisões críticas, limitações \\
50--77 & Cap. 4 --- Resultados & Três atos, quatro frentes e os sete resultados superados (§4.6) \\
78--86 & Cap. 5 --- Discussão & Posicionamento e o quadro do grau de estabelecimento por frente \\
87--88 & Cap. 6 --- Cronograma & Depósito em dezembro, defesa em janeiro \\
89--92 & Referências & 38 obras, todas lidas integralmente \\
94--107 & Apêndice A & As equações e \emph{todas} as especificações rodadas \\
108--112 & Apêndice B & As 31 decisões numeradas \\
113--115 & Apêndice C & As 11 rotinas citadas nominalmente \\
\bottomrule
\end{longtable}

\vspace{12pt}
Três características desse arranjo respondem sozinhas a perguntas frequentes de
banca, e vale conhecê-las como argumentos, e não apenas como fatos.

\vspace{4pt}
\textbf{O Apêndice A traz \emph{todas} as especificações rodadas, e não as que
convêm.} Ele é gerado por rotina a partir dos arquivos que os \emph{pipelines}
gravaram, de modo que nenhum coeficiente foi digitado à mão. Se a banca
perguntar por uma especificação que o Capítulo 4 não comenta, é grande a chance
de ela estar lá, inclusive as nulas.

\textbf{A bibliografia tem 38 entradas, e isso é uma escolha declarada.} A regra
adotada foi a de que só entra nas referências obra lida integralmente; o resto
está na §2.10 como lista de leitura, com forma bibliográfica completa e a
indicação do que se espera de cada uma. Uma bibliografia curta e conferida é
mais defensável que uma longa e citada de segunda mão, e a §2.10 mostra que o
levantamento foi feito.

\textbf{As decisões metodológicas são numeradas e consultáveis.} Quando o texto
diz ``decisão D26'', existe entrada correspondente no Apêndice B com o que se
decidiu e por quê. Isso permite responder a ``por que você fez assim?'' com uma
página, e não com uma reconstrução de memória.
